\section{A Comparative View: R1CS, PLONK, and Halo2 for Neural Networks}

\subsection{Running Example: A Small MLP}

Consider a simple neural network:
\[
z_1 = 2x_1 - x_2 + 1, \qquad
z_2 = x_1 + 3x_2,
\]
\[
h_1 = \operatorname{ReLU}(z_1), \qquad
h_2 = \operatorname{ReLU}(z_2),
\]
\[
y = 4h_1 - 2h_2 + 3.
\]

Our goal is to express this computation in three different proof systems.

\medskip

\noindent
\textbf{Key question.}  
How does the same computation look under different constraint paradigms?

\subsection{R1CS Representation}

R1CS expresses computation using constraints:
\[
\langle A_i, w \rangle \cdot \langle B_i, w \rangle = \langle C_i, w \rangle.
\]

Affine operations are encoded as:
\[
(\text{linear form}) \cdot 1 = \text{output}.
\]

For example:
\[
(1 + 2x_1 - x_2)\cdot 1 = z_1,
\qquad
(x_1 + 3x_2)\cdot 1 = z_2.
\]

The output layer becomes:
\[
(3 + 4h_1 - 2h_2)\cdot 1 = y.
\]

\paragraph{Handling ReLU in R1CS.}

R1CS does not support comparisons directly, so two approaches are used.

\medskip

\noindent
\textbf{(A) Polynomial approximation.}
\[
h = p(z), \quad p(z) = a_0 + a_1 z + a_2 z^2.
\]

This requires:
\[
z \cdot z = u, \qquad
(a_0 + a_1 z + a_2 u)\cdot 1 = h.
\]

\medskip

\noindent
\textbf{(B) Gadget-based emulation.}  
More complex constructions can simulate comparisons, but at significant cost.

\medskip

\noindent
\textbf{Observation.}  
R1CS favors algebraic (polynomial) representations over piecewise functions.

\subsection{PLONK Representation}

PLONK represents computation as a table of values with row-wise constraints:
\[
q_M ab + q_L a + q_R b + q_O c + q_C = 0.
\]

A possible layout is:

\begin{center}
\begin{tabular}{c|ccc|l}
Row & $a$ & $b$ & $c$ & Meaning \\
\hline
1 & $x_1$ & $x_2$ & $z_1$ & affine \\
2 & $x_1$ & $x_2$ & $z_2$ & affine \\
3 & $z_1$ & -- & $h_1$ & activation \\
4 & $z_2$ & -- & $h_2$ & activation \\
5 & $h_1$ & $h_2$ & $y$ & output
\end{tabular}
\end{center}

Affine rows are enforced by choosing coefficients appropriately, e.g.
\[
2a - b + 1 - c = 0.
\]

\paragraph{ReLU via lookup.}

Define a table:
\[
T = \{(t, \max(0,t)) : t \in D\}.
\]

Then enforce:
\[
(z_1, h_1) \in T, \qquad (z_2, h_2) \in T.
\]

\medskip

\noindent
\textbf{Observation.}  
PLONK naturally supports lookup constraints, making piecewise functions efficient.

\subsection{Halo2 Representation}

Halo2 is a PLONKish framework with flexible circuit design:

\begin{itemize}
\item advice columns (witness values)
\item fixed columns (constants)
\item selector columns (activate gates)
\item lookup tables (nonlinearities)
\end{itemize}

\paragraph{Affine gate.}

Define:
\[
s_{\mathrm{aff}} \cdot (\alpha a + \beta b + \gamma - c) = 0.
\]

Use it for:
\[
2x_1 - x_2 + 1 = z_1, \quad
x_1 + 3x_2 = z_2, \quad
4h_1 - 2h_2 + 3 = y.
\]

\paragraph{ReLU via lookup.}

Construct a lookup table:
\[
(t, \operatorname{ReLU}(t)).
\]

Then enforce:
\[
(z_1, h_1) \text{ in table}, \qquad
(z_2, h_2) \text{ in table}.
\]

\paragraph{Alternative: polynomial gate.}

Define:
\[
s_{\mathrm{poly}} \cdot \big(c - (a_0 + a_1 a + a_2 a^2)\big) = 0.
\]

\medskip

\noindent
\textbf{Observation.}  
Halo2 combines custom gates and lookup arguments in a programmable framework.

\subsection{Lookup vs Polynomial ReLU}

\paragraph{Lookup-based ReLU.}
\[
(z,h) \in \{(t,\max(0,t))\}.
\]

\begin{itemize}
\item exact on discrete domain
\item efficient in PLONK/Halo2
\item requires quantization
\end{itemize}

\paragraph{Polynomial ReLU.}
\[
h = p(z).
\]

\begin{itemize}
\item compatible with R1CS
\item smooth algebraic form
\item only approximate
\end{itemize}

\subsection{Comparison}

\begin{center}
\begin{tabular}{l|c|c|c}
Feature & R1CS & PLONK & Halo2 \\
\hline
Constraint form & $a\cdot b=c$ & polynomial gate & programmable gates \\
Custom operations & limited & yes & yes \\
Lookup support & weak & strong & very strong \\
ReLU implementation & polynomial & lookup & lookup (preferred) \\
Circuit flexibility & low & medium & high
\end{tabular}
\end{center}

\subsection{Key Insight}

\begin{quote}
The same neural network can be expressed at multiple levels of abstraction:
\[
\text{R1CS} \;\subset\; \text{PLONK} \;\subset\; \text{Halo2}.
\]
\end{quote}

\noindent
R1CS provides a minimal algebraic foundation,  
PLONK introduces structured polynomial constraints,  
and Halo2 enables practical, modular circuit design.

\medskip

\noindent
\textbf{Conclusion.}  
The choice of representation is not merely technical—it determines how efficiently machine learning models can be proven in zero knowledge.